\begin{abstract}
Spatio-temporal sensor data in real-world systems is often sparse, noisy, and irregular, making it difficult to infer global structure from limited observations. Under extreme sparsity, we run into the limits of identifiability of latent system states, making latent field reconstruction fundamentally underconstrained. In such scenarios, multiple physically plausible fields may remain consistent with the same observations, requiring reconstruction models to rely heavily on inductive biases regarding locality, transport structure, and spatial regularity. 

Under such sparsity regimes, reliable reconstruction becomes concentrated around the observational support induced by the sensor network, making sensor-space modeling a more identifiable objective than unconstrained global field recovery. We introduce FieldFormer, a mesh-free transformer architecture designed for locality-aware sensor-space modeling in persistent sensor networks. For each query, FieldFormer aggregates local evidence using a learnable velocity-scaled distance metric that adapts neighborhood geometry to heterogeneous spatio-temporal relationships. Neighborhoods are constructed as fixed maximal sparse contexts over nearby sensors and bounded temporal windows, while learned velocity-scaled offsets modulate token geometry within this context, enabling stable and scalable inference under extreme sparsity. A local transformer encoder integrates neighborhood information, while global consistency is modeled through coordinate-based neural field formulation.

We evaluate FieldFormer across five benchmarks spanning synthetic and real-world spatio-temporal systems, including anisotropic heat diffusion, shallow-water dynamics, atmospheric transport fields, and pollution monitoring datasets. Our results reveal that locality-aware reconstruction provides strong advantages in persistent sparse sensor networks where local domains of dependence remain observed, enabling FieldFormer to consistently outperform state-of-the-art baselines on sensor-space prediction tasks under highly sparse and noisy sensing regimes. \\

\noindent\textbf{Code Repository}
Our code repository is available at \url{https://github.com/ankitbha/fieldformer}. The README file contains detailed instructions regarding the code base.
\end{abstract}
